\documentclass[12pt]{article}
\usepackage{fontspec}
\usepackage[english, russian]{babel}
\usepackage{amsmath, amssymb}
\usepackage{tikz}
\usepackage{float}

\usepackage{pgfplots}
\pgfplotsset{compat=newest}
\usetikzlibrary{arrows.meta}

\setmainfont{CMU Serif}
\setsansfont{CMU Sans Serif}
\setmonofont{CMU Typewriter Text}

\usepackage{hyperref}
\hypersetup{
    colorlinks=true,        % цветные ссылки (false - черные рамки)
    linkcolor=black,        % цвет ссылок на разделы
    urlcolor=blue,          % цвет URL-ссылок
    citecolor=green,        % цвет ссылок на литературу
    pdfborder={0 0 0}       % убрать рамки вокруг ссылок
}

\usepackage[a4paper,
            left=3cm,
            right=3cm,
            top=2cm,
            bottom=3cm]{geometry}

\begin{document}

\begin{titlepage}
    \begin{center}
        \large
        \textsc{ФГАОУ ВО «НИУ ИТМО»\\
        Факультет программной инженерии и~компьютерной~техники}
        \vspace*{4cm}
            
        \LARGE
        \textbf{Лабораторная работа №4}

        \vspace{0.5cm}
        \LARGE

        Численные методы безусловной оптимизации функции нескольких переменных
            
        \vspace{0.5cm}
        \large
        (по дисциплине «Методы оптимизации»)

        \vspace{1.5cm}

        \hfill\large
        \begin{tabular}{@{}l@{}}
        \textbf{Выполнили:}\\[2mm]
        Лабин Макар Андреевич,\\
        МОПТ 3.3, P3231, 466449;\\
        Ожеховский Александр Сергеевич,\\
        МОПТ 3.3, P3231, 466944.\\[2mm]
        \textbf{Проверил:}\\[2mm]
        Запорожцев Иван Федорович
        \end{tabular}
            
        \vfill
            
        \vspace{0.8cm}
            
        \includegraphics[width=0.4\textwidth]{itmo}
            
        \Large
        г. Санкт-Петербург, Россия\\
        2026
            
    \end{center}
\end{titlepage}

\renewcommand{\thefootnote}{[\arabic{footnote}]}

\tableofcontents
\newpage

\section{Постановка задач}

В рамках выполнения лабораторной работы \textit{по варианту №[[VARIANT]]} необходимо исследовать следующую функцию на наличие экстремумов:
\[
  f(x, y) = [[A]]x^2 + [[B]]y^2 + [[C]]xy + [[D]]x + [[E]]y + [[F]]
\]

Исследования необходимо производить двумя способами:
\begin{itemize}
  \item аналитически;
  \item численно:
  \begin{itemize}
    \item методом циклического покоординатного спуска;
    \item методом градиентного спуска;
    \item методом наискорейшего спуска.
  \end{itemize}
\end{itemize}

Для численных методов необходимо использовать следующие параметры:
\begin{itemize}
  \item Начальное приближение: \( x^{(0)}=[[X0]] \).
  \item Для критерия останова: \( \varepsilon = [[EPSILON]] \).
\end{itemize}

\newpage

\section{Аналитический метод}

Данная функция \( f(x,y) \) --- квадратичная форма. Для поиска её экстремумов воспользуемся необходимым и достаточным условиями экстремума функции нескольких переменных.

\subsection{Необходимое условие экстремума}

Найдём частные производные первого порядка и приравняем их к нулю, чтобы определить координаты стационарных точек:
\[
  \begin{cases}
    \frac{\partial f}{\partial x} = 2 \cdot [[A]]x + [[C]]y + [[D]] = 0 \\
    \frac{\partial f}{\partial y} = 2 \cdot [[B]]y + [[C]]x + [[E]] = 0
  \end{cases}
\]

Решая данную систему, находим координаты единственной стационарной точки:
\[
x^* = [[X_STAR]], \quad y^* = [[Y_STAR]]
\]

\subsection{Достаточное условие экстремума}

Для проверки характера найденной точки составим матрицу Гессе:
\begin{equation*}
H(x, y) = \begin{pmatrix}
    \frac{\partial^2 f}{\partial x^2} & \frac{\partial^2 f}{\partial x \partial y} \\[1ex]
    \frac{\partial^2 f}{\partial y \partial x} & \frac{\partial^2 f}{\partial y^2}
\end{pmatrix}
=
\begin{pmatrix}
    2 \cdot [[A]] & [[C]] \\
    [[C]] & 2 \cdot [[B]]
\end{pmatrix}
\end{equation*}

Исследуем знаки главных миноров гессиана:
\begin{itemize}
    \item \( \Delta_1 = \frac{\partial^2 f}{\partial x^2} = 2 \cdot [[A]] = [[DELTA_1]]>0 \).
    \item \( \Delta_2 = \det(H) = 4 \cdot [[A]] \cdot [[B]] - [[C]]^2 = [[DELTA_2]]>0 \).
\end{itemize}

По критерию Сильвестра, матрица Гессе положительно определена, то есть \( f(x,y) \) строго выпукла вниз.

Следовательно, стационарная точка есть \textit{точка глобального минимума}:
\[
x^* = ([[X_STAR]], [[Y_STAR]])^T
\]

Значение функции в точке минимума:
\[
f(x^*, y^*) = [[F_STAR]]
\]

\newpage
\section{Метод покоординатного спуска}

\subsection{Применимость метода}

Целевая функция ограничена снизу, что гарантирует сходимость метода. Задачи одномерной минимизации решаются аналитически за константное время путём поиска вершины параболы для данного направления.

\subsection{Подробный расчет итераций}
[[CYCLIC_ITERATIONS]]

\subsection{Сводная таблица итераций}
\begin{table}[H]
\centering
\begin{tabular}{|c|c|c|c|c|}
\hline
$k$ & $x^{(k-1)}$ & $\alpha^{(k)}$ & $x^{(k)}$ & $|x^{(k)} - x^{(k-1)}|$ \\
\hline
[[CYCLIC_TABLE]]
\end{tabular}
\end{table}

Итого имеем:
\[
  \tilde{x}^{\ast}=[[FINAL_CYCLIC]]
\]

\subsection{График спусков}

\begin{figure}[H]
  \centering
  \begin{tikzpicture}
    \begin{axis}[
      width=10cm, height=10cm,
      axis lines=middle,
      xlabel={$x$},
      ylabel={$y$},
      xlabel style={right},
      ylabel style={above},
      colormap={blackwhite}{gray(0cm)=(0.4); gray(1cm)=(0.9)},
      view={0}{90}
    ]
      \addplot3[
        contour lua={
          number=10,
          labels=false,
        },
        samples=30,
        domain=[[X_MIN_CYCLIC]]:[[X_MAX_CYCLIC]],
        domain y=[[Y_MIN_CYCLIC]]:[[Y_MAX_CYCLIC]],
      ] {[[A]]*x^2 + [[B]]*y^2 + [[C]]*x*y + [[D]]*x + [[E]]*y + [[F]]};

      [[CYCLIC_PATH]]
        
    \end{axis}
  \end{tikzpicture}
  \caption{Все итерации метода циклического покоординатного спуска.}
\end{figure}

\newpage

\section{Метод градиентного спуска}

\subsection{Применимость метода}

Целевая функция ограничена снизу, что гарантирует сходимость метода. Также она непрерывно дифференцируема на \( \mathbb{R}^2 \), что позволяет считать антиградиенты в любой точке плоскости.

Установим скорость обучения постоянной, независимой от номера итерации:
\[
  \eta = [[GD_ETA]]
\]

\subsection{Подробный расчет итераций}
[[GD_ITERATIONS]]

\subsection{Сводная таблица итераций}
\begin{table}[H]
\centering
\begin{tabular}{|c|c|c|c|c|}
\hline
$k$ & $x^{(k-1)}$ & $\omega^{(k)}$ & $x^{(k)}$ & $|\omega^{(k)}|$ \\
\hline
[[GD_TABLE]]
\end{tabular}
\end{table}

Итого имеем:
\[
  \tilde{x}^{\ast}=[[FINAL_GD]]
\]

\subsection{График спусков}

\begin{figure}[H]
  \centering
  \begin{tikzpicture}
    \begin{axis}[
      width=10cm, height=10cm,
      axis lines=middle,
      xlabel={$x$},
      ylabel={$y$},
      xlabel style={right},
      ylabel style={above},
      colormap={blackwhite}{gray(0cm)=(0.4); gray(1cm)=(0.9)},
      view={0}{90}
    ]
      \addplot3[
        contour lua={
          number=10,
          labels=false,
        },
        samples=30,
        domain=[[X_MIN_GD]]:[[X_MAX_GD]],
        domain y=[[Y_MIN_GD]]:[[Y_MAX_GD]],
      ] {[[A]]*x^2 + [[B]]*y^2 + [[C]]*x*y + [[D]]*x + [[E]]*y + [[F]]};

      [[GD_PATH]]
        
    \end{axis}
  \end{tikzpicture}
  \caption{Все итерации метода градиентного спуска.}
  \end{figure}

\newpage

\section{Метод наискорейшего спуска}

\subsection{Применимость метода}

Целевая функция ограничена снизу, что гарантирует сходимость метода. Также она непрерывно дифференцируема на \( \mathbb{R}^2 \), что позволяет считать антиградиенты в любой точке плоскости. Задачи одномерной минимизации решаются аналитически за константное время путём поиска вершины параболы для данного направления.

\subsection{Подробный расчет итераций}
[[STEEPEST_ITERATIONS]]

\subsection{Сводная таблица итераций}
\begin{table}[H]
\centering
\begin{tabular}{|c|c|c|c|c|c|}
\hline
$k$ & $x^{(k-1)}$ & $u^{(k)}$ & $\beta^{(k)}$ & $x^{(k)}$ & $|\omega^{(k)}|$ \\
\hline
[[STEEPEST_TABLE]]
\end{tabular}
\end{table}

Итого имеем:
\[
  \tilde{x}^{\ast}=[[FINAL_STEEPEST]]
\]

\subsection{График спусков}

\begin{figure}[H]
  \centering
  \begin{tikzpicture}
    \begin{axis}[
      width=10cm, height=10cm,
      axis lines=middle,
      xlabel={$x$},
      ylabel={$y$},
      xlabel style={right},
      ylabel style={above},
      colormap={blackwhite}{gray(0cm)=(0.4); gray(1cm)=(0.9)},
      view={0}{90}
    ]
      \addplot3[
        contour lua={
          number=10,
          labels=false,
        },
        samples=30,
        domain=[[X_MIN_STEEPEST]]:[[X_MAX_STEEPEST]],
        domain y=[[Y_MIN_STEEPEST]]:[[Y_MAX_STEEPEST]],
      ] {[[A]]*x^2 + [[B]]*y^2 + [[C]]*x*y + [[D]]*x + [[E]]*y + [[F]]};

      [[STEEPEST_PATH]]
        
    \end{axis}
  \end{tikzpicture}
  \caption{Все итерации метода наискорейшего спуска.}
  \end{figure}

\end{document}